% CVPR two-column report. Put the official CVPR author-kit files (cvpr.sty, ieeenat_fullname.bst)
% from https://github.com/cvpr-org/author-kit next to this file (or use the CVPR template on Overleaf).
% Without them the file still compiles with a plain two-column layout.
\documentclass[10pt,twocolumn,letterpaper]{article}
\IfFileExists{cvpr.sty}{\usepackage[final]{cvpr}}{\usepackage[margin=0.75in]{geometry}}
\usepackage{graphicx,booktabs,amsmath,amssymb,xcolor,url,adjustbox}
\usepackage[hidelinks]{hyperref}
\graphicspath{{../figures/}}
\newcommand{\todo}[1]{\textcolor{red}{[TODO: #1]}}
\newcommand{\tab}[1]{\IfFileExists{tables/#1.tex}{\input{tables/#1.tex}}{\todo{run make\_tables.py: #1}}}
\newcommand{\figw}[2]{\begin{figure*}[t]\centering\IfFileExists{../figures/#1}{\includegraphics[width=0.85\textwidth]{#1}}{\todo{figure #1}}\caption{#2}\label{fig:#1}\end{figure*}}
\newcommand{\fig}[3]{\begin{figure}[t]\centering\IfFileExists{../figures/#1}{\includegraphics[width=#2]{#1}}{\todo{figure #1}}\caption{#3}\label{fig:#1}\end{figure}}

\title{MLPs, CNNs, Autoencoders, SVMs and Clustering on CIFAR10}
\author{Bighnaraj Mohapatra \quad Roll no.\ M26CSA011\\
Repository: \url{https://github.com/TODO/a3-cifar10} \quad Final commit: \texttt{TODO-commit-hash}}

\begin{document}
\maketitle

\begin{abstract}
\todo{Write last, 4--5 sentences, one finding per question, each with its measured number and the
comparison it comes from. Do not state anything the tables do not show.}
\end{abstract}

%=====================================================================
\section{Common setup}
\label{sec:setup}
From the 50,000 CIFAR10~\cite{krizhevsky2009cifar} training images we draw one stratified split of 45,000 training and 5,000
validation images (500 per class) with \texttt{train\_test\_split(stratify=y, random\_state=26011)};
26011 is the digit sequence of the roll number. The indices are saved in
\texttt{splits/split\_seed26011.npz} and every experiment loads that file. Validation data are used only
for model, architecture and hyper-parameter selection. The official test set is evaluated by a separate
script only after all choices are fixed.

Images are scaled to $[0,1]$ and standardised per channel with statistics of the (transformed) training
images. No data augmentation is used: random crops and flips encode a spatial prior, which would confound
the MLP--CNN comparisons of Q1 and Q2 and has no meaning for permuted images. All classifiers are trained
with AdamW (weight decay $5\cdot10^{-4}$), batch size 128 and a cosine schedule for \todo{60} epochs. The
learning rate of every compared model is chosen from the same grid $\{3\cdot10^{-4},10^{-3},3\cdot10^{-3}\}$
by validation accuracy (seed 26011). The checkpoint with the best validation accuracy is kept. Final
comparisons use seeds 26011--26013 and report mean $\pm$ sample standard deviation. Fairness criteria and
predictions were committed in \texttt{PREREGISTRATION.md} (commit \todo{hash}) before the final runs.
\textbf{Reuse.} Components are introduced once and reused: MLP-B and CNN-A (Q1(a)) are the models of
Q2(a), and CNN-A is the fixed classifier of Q3(b). The Q3(a) autoencoders are the baseline of Q3(b), are
probed in Q3(c) and provide the Q4 representation (Table~\ref{tab:models} and the README reuse map).

%=====================================================================
\section{Q1: Why does an MLP struggle on CIFAR10?}
\subsection{Q1(a): separating the reason from an alternative}
\textbf{Chosen reason (H1).} The MLP has no spatial inductive bias. It sees an unordered 3,072-vector,
with neither locality nor weight sharing, so it must learn pixel relations from data and generalises
poorly even when it can fit the training set. \textbf{Alternative (H2).} The MLP under-fits because of
too little capacity or poor optimisation. The two hypotheses make different predictions. H2 predicts low
training accuracy, validation accuracy that rises with width, and large gains from better optimisation
(BN, tuned LR). H1 predicts high training accuracy with a large train--val gap, validation accuracy that
saturates as width grows, and a clear gain from a spatial prior at equal parameter count.
\textbf{Design.} (i) A width sweep MLP-S/B/L/XL (0.41M--4.21M parameters). (ii) An optimisation control,
MLP-B-plain (no BN, untuned LR). (iii) A prior-injection control, CNN-A, parameter-matched to MLP-B
(816,202 vs 856,074, within 5\%). All use the protocol of Sec.~\ref{sec:setup}. Compute is not matched
(Table~\ref{tab:models}); we report it rather than hide it.

\tab{tab_models}
\tab{tab_q1a}
\figw{q1a_curves.png}{Q1(a)/Q2(a): training and validation accuracy per epoch (seed 26011).}
\fig{q1a_width_sweep.png}{0.85\linewidth}{Q1(a): MLP width sweep vs the parameter-matched CNN-A.}

\textbf{Interpretation.} \todo{Use Table~\ref{tab:q1a} and the figures: (1) is MLP-B's training accuracy
high and its train--val gap large? (2) does validation accuracy saturate with width? (3) how much did BN and
LR tuning add (MLP-B vs MLP-B-plain)? (4) CNN-A vs MLP-B at matched parameters. Conclude for H1 or H2 only
as far as these numbers allow.}

\subsection{Q1(b): is structure learned before memorisation?}
\textbf{Setting.} MLP-Mem (3072--2048--2048--10, 10,510,346 parameters, no BN, no weight decay, no
dropout) is trained with Adam (constant LR $3\cdot10^{-4}$) for \todo{200} epochs. A fixed 40\% of the
training labels (seed 26011) is replaced by a label drawn uniformly from the \emph{other} nine classes, so
fitting a corrupted label is unambiguously memorisation. Validation labels stay clean. Every epoch we
measure accuracy on the clean-label subset, accuracy on the noisy subset against the assigned (wrong) and
the true labels, and validation accuracy. Two controls use 0\% and 100\% corrupted labels.
\textbf{Pre-registered rule.} Let $t^*$ be the epoch of peak validation accuracy. We conclude that
structure is learned first iff, at $t^*$, the noisy-subset fit to assigned labels is below 20\% and the
noisy-subset accuracy against true labels is above chance (10\%).

\figw{q1b_memorisation.png}{Q1(b): training trajectory under 40\% label noise (left) and the
speed of fitting real vs random labels (right).}
\tab{tab_q1b}

\textbf{Interpretation.} \todo{Apply the rule using Table~\ref{tab:q1b}. Compare the first epochs at which
the clean and noisy subsets reach 25/50/90\%. Describe what happens to validation and noisy-vs-true accuracy
after $t^*$. Report test accuracy at $t^*$ vs the final epoch.}

%=====================================================================
\section{Q2: What does the CNN use?}
\subsection{Q2(a): a fair comparison and a proposed reason}
We reuse MLP-B and CNN-A from Q1(a). The fairness criterion was fixed before the final results
(\texttt{PREREGISTRATION.md}): parameters within $\pm10\%$, identical data, preprocessing, optimiser,
epochs, LR grid and checkpoint rule, no augmentation, and three seeds. \textbf{Proposed reason:} CNN-A's
advantage comes from \emph{locality and weight sharing}. Its small filters apply the same local detectors
at every position and so exploit correlations between neighbouring pixels. \todo{Support with the
train--val gaps and per-class accuracies from Q1(a).}

\subsection{Q2(b): same pixel values, disturbed spatial structure}
We apply one fixed random permutation of the 1,024 pixel positions (\texttt{default\_rng(26011)}),
identical for all channels and for every training, validation and test image, and retrain both models
with the same hyper-parameters. \textbf{Prediction stated beforehand.} MLP-B is unchanged within seed
noise: a fixed input permutation is equivalent to permuting its first-layer weights, whose i.i.d.\
initialisation is permutation-invariant. CNN-A drops substantially. If its advantage disappears, the
locality explanation of Q2(a) is supported. If it stays clearly ahead, it is not.

\tab{tab_q2}

\subsection{Q2(c): colour statistics or luminance structure?}
\textbf{Explanations.} E1: CNN-A relies mainly on colour statistics. E2: it relies mainly on luminance
spatial structure. \textbf{Experiment.} Both models are retrained and tested on grayscale images
($Y=0.299R+0.587G+0.114B$, replicated to three channels, so architecture and parameters are unchanged).
As a separately labelled diagnostic, the RGB-trained models are also evaluated on grayscale test images.
\textbf{Errors and confidence} (test set): per-class accuracy changes, new confusions, mean max-softmax for
correct and incorrect predictions, ECE with 15 bins~\cite{guo2017calibration}, and the share of errors made
with confidence $>0.9$ (\texttt{results/q2c\_confidence.json}).

\fig{q2c_reliability.png}{0.8\linewidth}{Q2(c): reliability diagrams for CNN-A (test set).}

\textbf{Conclusion linking (a)--(c).} \todo{3--4 sentences: the reason proposed in (a); whether the
permutation result in (b) supports it, with the numbers; what (c) says about colour vs spatial structure.}

%=====================================================================
\section{Q3: Autoencoders}
\subsection{Q3(a): bottleneck capacity}
The three fully convolutional AEs share every layer except the final $1\times1$ encoder and first
$1\times1$ decoder convolutions. The bottleneck is $z\in\mathbb{R}^{c\times4\times4}$ with
$c\in\{2,16,128\}$, i.e.\ $d\in\{32,256,2048\}$ (compression $96\times$, $12\times$, $1.5\times$). Training
is identical: per-pixel MSE on $[0,1]$, Adam with LR $10^{-3}$, cosine schedule, \todo{50} epochs, batch
128, and the checkpoint with the lowest validation MSE. Selection uses validation MSE only.
\tab{tab_q3a}
\textbf{Choice.} \todo{Name the selected $d$, justify it with validation MSE and the train/val gap
(Fig.~\ref{fig:q3a_mse_curves.png}), and give its test MSE.}
\fig{q3a_mse_curves.png}{0.85\linewidth}{Q3(a): reconstruction MSE per epoch.}

\subsection{Q3(b): denoising}
\textbf{Corruption.} Additive Gaussian noise, $\sigma=0.2$ on $[0,1]$, clipped. It is a corruption a
reconstruction model can plausibly undo, because i.i.d.\ noise lies off the image
manifold~\cite{vincent2008denoising}. $\sigma$ was fixed in advance. A validation scan
(\texttt{results/q3b\_sigma\_scan\_val.json}) confirms that it hurts the classifier. \textbf{Protocol.}
The DAE has the same architecture as the selected clean AE and is trained on noisy inputs with clean
targets, without labels. Fresh noise is drawn each batch. Validation and test use fixed seeded noise.
The fixed classifier is the clean CNN-A (Q2(a), seed 26011). AE(clean) is a control for accuracy lost to
reconstruction blur alone.
\tab{tab_q3b}
\figw{q3b_examples.png}{Q3(b): test images: clean, noisy, AE(noisy), DAE(noisy).}
\textbf{Decision.} \todo{Compare DAE(noisy) with noisy (did it help?) and with AE(clean) and clean (how much
of the recoverable accuracy is recovered?); use the agreement column and the per-class results.}

\subsection{Q3(c): class information vs reconstruction}
Encoders are frozen. For each representation (the three AEs, the DAE, an untrained encoder of the selected
architecture, and raw pixels) we fit a multinomial logistic-regression probe (L2 strength chosen on
validation) and a kNN probe ($k$ chosen on validation) on the 45k training features, then report test
accuracy. Probes after PCA to 32 dimensions control for the different dimensionalities.
\tab{tab_q3c}
\textbf{Do the two quantities agree?} \todo{Compare the ranking by test MSE with the ranking by probe
accuracy. With four trained models, describe the pattern rather than claim a correlation.}

%=====================================================================
\section{Q4 (bonus): SVMs and clustering}
\textbf{Fixed representation.} The $d=256$ bottleneck of the Q3(a) AE, standardised. It was learned
without labels, so training-set features carry no label memorisation. \textbf{Fixed subset.} 10,000
training images (1,000 per class, seed 26011), chosen before any model selection and used for every SVM
and clustering fit. Validation uses the full 5k and test the full 10k.

\subsection{Q4(a): $C$ and kernel}
Kernels: linear, RBF and polynomial (degree 3, $c_0=1$), with $C\in\{0.01,0.1,1,10,100\}$ and $\gamma$
fixed at the numeric value of \texttt{'scale'}. Selection uses validation accuracy only; the selected SVM is
evaluated once on test.
\tab{tab_q4a}
\figw{q4a_sweep.png}{Q4(a): validation/training accuracy and support-vector fraction vs $C$.}
\textbf{Under- and over-fitting.} \todo{Use the train--val gap and the SV fraction across $C$, per kernel.}

\subsection{Q4(b): difficult pairs}
The three pairs with the largest symmetric test confusion are analysed with the pairwise SVM (same kernel,
$C$ and numeric $\gamma$; exactly the one-vs-one classifier of the selected SVM).
\tab{fig_q4b}
\todo{For each pair: share of test errors inside the margin band $|f|<1$ vs confident errors, bounded vs
free SVs, and the Fisher ratio from \texttt{results/q4b\_pairs.json}; relate them to the images.}

\subsection{Q4(c): clustering vs classification}
$k$-means and diagonal GMMs with $k\in\{5,8,10,12,15,20\}$ are fitted on the subset without labels. The
method and $k$ are chosen by validation silhouette. Test representations are assigned without labels, and
labels are used only to evaluate.
\tab{tab_q4c}
\fig{q4c_pairs.png}{0.8\linewidth}{Q4(c): per class pair, SVM confusion rate vs clustering ARI (test).}
\textbf{Disagreement.} \todo{Name one class or pair where the SVM and the clustering disagree, quoting the
SVM pair confusion rate and the cluster pair ARI from \texttt{results/q4c\_pairs\_test.csv}.}

{\small
\IfFileExists{ieeenat_fullname.bst}{\bibliographystyle{ieeenat_fullname}}{\bibliographystyle{plain}}
\bibliography{references}
}
\end{document}
